\chapter{Kiến trúc hệ thống}
\label{chap:system-architecture}

\section{Kiến trúc logic}

Kiến trúc được chia thành năm lớp: thiết bị/publisher, messaging, thu nhận ứng dụng, persistence và quản trị. Việc phân lớp này bám theo các entry point \texttt{SimulatorApp::loop}, mqtt.js \texttt{connectToHiveMQ}, Slim \path|public/index.php|, migration dùng chung và Laravel \path|routes/web.php| (E-SIM-02, E-DM-02, E-DM-06, E-DB-01, E-ADM-04).

\begin{landscape}
\begin{figure}[p]
\centering
\resizebox{0.96\linewidth}{!}{%
\begin{tikzpicture}[
  node distance=10mm,
  comp/.style={draw=SmartWaterBlue,rounded corners=3pt,align=center,text width=3.35cm,minimum height=1.55cm,fill=SmartWaterBlue!7,font=\small},
  data/.style={comp,draw=StatusImplemented,fill=StatusImplemented!8},
  legacy/.style={comp,draw=StatusPlanned,dashed,fill=black!6},
  flow/.style={-{Latex[length=2.5mm]},thick,draw=SmartWaterNavy},
  dashedflow/.style={flow,dashed,draw=StatusPartial}
]
\node[comp] (mqtt) {mqtt.js Client\\Nhận MQTT/WSS\\E-DM-06};
\node[comp,right=of mqtt] (route) {Slim Routes\\\texttt{public/index.php}\\E-DM-02};
\node[comp,right=of route] (service) {Telemetry Service\\Normalize payload\\E-DM-03};
\node[comp,right=of service] (repo) {Telemetry Repository\\Insert / latest / summary\\E-DM-04};
\node[data,right=of repo] (telemetry) {MySQL canonical\\Bảng \texttt{telemetry}\\E-DB-02};

\node[comp,below=24mm of telemetry] (models) {Laravel Models\\Telemetry / MCU / Device\\E-ADM-05};
\node[comp,left=of models] (view) {Admin Controller + Blade\\TDS / cảnh báo\\Partially Implemented};
\node[legacy,below=24mm of mqtt] (worker) {.NET Worker\\DeviceService + Repository\\Contract legacy};
\node[legacy,right=of worker] (legacydb) {MySQL legacy\\\texttt{devices, device\_data}\\E-MQ-05};

\draw[flow] (mqtt) -- node[above]{HTTP} (route);
\draw[flow] (route) -- (service);
\draw[flow] (service) -- (repo);
\draw[flow] (repo) -- node[above]{PDO} (telemetry);
\draw[flow] (telemetry) -- node[right]{Eloquent} (models);
\draw[flow] (models) -- node[above]{query} (view);
\draw[dashedflow] (worker) -- node[above]{Dapper} (legacydb);
\end{tikzpicture}%
}

\vspace{3mm}
\small\textbf{Chú giải:} nhóm xanh biểu diễn pipeline canonical; nhóm xám nét đứt
là worker và schema legacy, không nối trực tiếp với bảng telemetry canonical.
\caption{Thiết kế component xử lý và tiêu thụ telemetry}
\label{fig:component-architecture}
\end{figure}
\end{landscape}


\subsection{Lớp publisher}

Simulator và firmware chính đều là publisher của contract chuẩn: khi MQTT đã kết nối, chúng phát JSON có \texttt{mcu\_id} dạng chuỗi và object \texttt{telemetry} lên \path|devices/telemetry|. Firmware chính tạo TDS và cảnh báo từ \texttt{SensorData}; simulator tạo TDS mẫu. Cả hai dùng định danh \texttt{MCU\_ID} cùng mô hình với schema canonical (E-FW-03, E-FW-04, E-SIM-03, E-SIM-04).

\subsection{Lớp messaging}

HiveMQ là phụ thuộc bên ngoài. Client ESP32 kết nối MQTT qua TLS; trình duyệt Device Monitor kết nối WSS bằng mqtt.js; worker .NET kết nối TCP/TLS bằng MQTTnet. Repository chỉ chứa key cấu hình client, không chứa Docker Compose, chart hay cấu hình triển khai broker, nên trạng thái được phân loại là \textit{Configured}, không phải broker do dự án triển khai (E-FW-04, E-DM-06, E-MQ-03, E-MQ-06).

\subsection{Lớp thu nhận và persistence}

Trong đường Device Monitor, trình duyệt vừa là MQTT subscriber vừa là HTTP adapter: callback \texttt{client.on('message')} giải mã JSON nếu có thể và gọi \texttt{sendTelemetry}; Slim API chuẩn hóa payload qua \texttt{TelemetryService} rồi \texttt{TelemetryRepository} thực hiện lệnh \texttt{INSERT INTO telemetry} (E-DM-02--E-DM-04, E-DM-06). \texttt{Database::pdo} tự mở PDO từ các key môi trường và kiểm tra bảng telemetry do migration dùng chung quản lý (E-DM-05).

Laravel có cổng ingest thứ hai tại \path|POST /api/telemetry|. \texttt{TelemetryController::ingest} chấp nhận payload phẳng hoặc object \texttt{payload}, kiểm tra \texttt{mcu\_id} là chuỗi tối đa 50 ký tự rồi tạo model \texttt{Telemetry}; file \path|app/Services/TelemetryService.php| không có implementation và không tham gia luồng này (E-ADM-05, E-ADM-06).

\subsection{Lớp quản trị}

Admin autoload factory và seeder từ dự án \texttt{smartwater-database}; migration dùng chung được chạy bằng \path|smartwater-database/migrate.bat| (E-ADM-03). 

Khi mở route \texttt{devices.show}, \texttt{DeviceController::show} gọi quan hệ \texttt{Device::telemetry}; quan hệ đi qua MCU gắn với product và view dựng đồ thị TDS cùng danh sách cảnh báo. Bằng chứng nằm tại \path|app/Models/Device.php|, \path|app/Models/Mcu.php|, \path|app/Http/Controllers/DeviceController.php| và \path|resources/views/devices/show.blade.php|.

\section{Kiến trúc triển khai}

\begin{landscape}
\begin{figure}[p]
\centering
\resizebox{0.92\linewidth}{!}{%
\begin{tikzpicture}[
  node distance=15mm and 22mm,
  host/.style={draw=SmartWaterBlue,rounded corners=3pt,align=center,text width=4.0cm,minimum height=1.65cm,fill=SmartWaterBlue!7,font=\small},
  external/.style={host,draw=StatusPartial,dashed,fill=StatusPartial!10},
  legacy/.style={host,draw=StatusPlanned,dashed,fill=black!6},
  flow/.style={-{Latex[length=2.5mm]},thick,draw=SmartWaterNavy},
  branch/.style={flow,dashed,draw=StatusPartial}
]
\node[host] (esp) {Edge Host\\ESP32 firmware / simulator\\PlatformIO};
\node[external,right=of esp] (hive) {External Service\\HiveMQ Broker\\Server manifest Not Found};
\node[host,right=of hive] (browser) {Operator Browser\\mqtt.js qua WSS\\Kết nối thủ công};

\node[legacy,below=of esp] (worker) {Windows / .NET Host\\MQTT Worker legacy\\Windows Service configured};
\node[host,below=of hive] (admin) {Admin Application Host\\PHP 8.3 / Laravel 13\\Web + API};
\node[host,below=of browser] (monitor) {Monitor Application Host\\PHP 8.1 / Slim 4\\HTTP API};

\node[legacy,below=of worker] (legacydb) {MySQL Legacy Schema\\\texttt{devices, device\_data}\\Dapper};
\node[host,below=of admin] (mysql) {MySQL Canonical Schema\\Shared migrations\\\texttt{mcus, telemetry}};

\draw[flow] (esp) -- node[above]{TLS} (hive);
\draw[flow] (hive) -- node[above]{WSS} (browser);
\draw[flow] (browser) -- node[right]{HTTP/JSON} (monitor);
\draw[flow] (monitor.south) |- node[pos=0.7,below]{PDO} (mysql.east);
\draw[flow] (admin) -- node[right]{Eloquent} (mysql);
\draw[branch] (hive.south) -- ++(0,-7mm) -| node[pos=0.65,above]{MQTT} (worker.north);
\draw[branch] (worker) -- node[left]{Dapper} (legacydb);
\end{tikzpicture}%
}

\vspace{3mm}
\small\textbf{Chú giải:} mỗi hộp là một ranh giới triển khai. Nét đứt chỉ thành phần
bên ngoài hoặc legacy; không biểu thị chúng đã được xác minh chạy đồng thời.
\caption{Topology triển khai suy ra từ manifest và entry point hiện có}
\label{fig:deployment-architecture}
\end{figure}
\end{landscape}


Các ranh giới chạy độc lập gồm hai project PlatformIO, Slim Device Monitor, Laravel Admin và .NET Worker. Database migration không phải ứng dụng chạy riêng; Admin đổi database path sang thư mục migration dùng chung. Broker nằm ngoài repository (E-REP-01, E-ADM-03, E-MQ-06; Phụ lục~\ref{app:repository-structure}).

\section{Giao diện giữa các thành phần}

\begin{table}[htbp]
\centering
\small
\begin{tabular}{p{0.21\textwidth}p{0.17\textwidth}p{0.26\textwidth}p{0.27\textwidth}}
\hline
\textbf{Nguồn--đích} & \textbf{Giao thức} & \textbf{Hợp đồng} & \textbf{Bằng chứng} \\
\hline
Simulator--HiveMQ & MQTT/TLS & Topic \path|TOPIC_PREFIX/telemetry|; JSON \texttt{mcu\_id} và \texttt{telemetry.tds} & \texttt{SimulatorApp.cpp}; \texttt{MqttManager.cpp} (E-SIM-03, E-SIM-04). \\
HiveMQ--trình duyệt & MQTT/WSS & Danh sách subscription từ cấu hình trình duyệt & \path|public/assets/app.js::connectToHiveMQ| (E-DM-06). \\
Trình duyệt--Slim & HTTP JSON & \path|POST /api/telemetry|: \texttt{topic, payload, source} & \texttt{app.js::sendTelemetry}; \path|public/index.php| (E-DM-02, E-DM-06). \\
Slim--MySQL & PDO & \path|INSERT/SELECT telemetry| & \texttt{TelemetryRepository.php}; \texttt{Database.php} (E-DM-04, E-DM-05). \\
Admin--MySQL & Eloquent & \texttt{Telemetry}, \texttt{Mcu}, quan hệ theo \texttt{mcu\_id} & \texttt{Telemetry.php}; \texttt{Device.php}; migration telemetry (E-ADM-05, E-DB-02). \\
HiveMQ--worker .NET & MQTT/TLS & DTO \texttt{device\_id}; ghi \texttt{device\_data} & \texttt{MqttService.cs}; \texttt{DeviceService.cs}; \texttt{DeviceRepository.cs} (E-MQ-03--E-MQ-05). \\
\hline
\end{tabular}
\caption{Giao diện tích hợp có bằng chứng trong source}
\label{tab:architecture-interfaces}
\end{table}

\section{Ranh giới không đồng bộ}

Firmware chính, simulator, Device Monitor, Laravel API và migration dùng đường chuẩn \path|mcu_id/telemetry|. Nhánh legacy còn lại là worker .NET: worker giải mã \texttt{device\_id} rồi ghi schema riêng \path|device_data|. Không có adapter source nào ánh xạ nhánh worker này sang pipeline canonical (E-FW-04, E-SIM-03, E-DM-03, E-ADM-05, E-DB-02, E-MQ-03--E-MQ-05).
